\documentclass{article}
\usepackage{../../Self_Style}

\title{Phys 20BL Week 2 Out of Lab: Electron Charge to Mass Ratio}
\author{Zih-Yu Hsieh}
\date{\today}

\begin{document}
\maketitle

\tableofcontents

\newpage

\section{Electron Energy}
\subsection{Exercise 11}
The kinetic energy of the electron that's emerging from the slit in the anode is given by the following:
\begin{equation}
    \frac{1}{2}mv^2 = e V
\end{equation}
In this case, $m$ represents the mass of the electron, $v$ represents the speed of the electron, $e$ represents the charge of the electron, and $V$ represents the electric potential of the environment (i.e. the electric potential within the anode in this case).

\hfil

\subsection{Exercise 12}

This equation does make sense to myself. Before the electrons were shot out from the anode (or, before the anode is provided any potential), one can assume the electrons on has nearly no energy (since there's no provided electric potential, so no potential energy; on the other hand, one can also assume the initial electrons are having low enough kinetic energy).

Which, after the anode is provided with some electric potential $V$, the electron's total energy is (about the same) as the potential energy $E = U = e V$; however, once the electron is shot out (no longer in the anode), one has $V_f = 0$ (the environmental potential can be treated as $0$), so $U_f = e V_f = 0$, showing that the total energy should nearly all transferred to kinetic energy. 

Hence, this resulted in $\frac{1}{2}mv^2 = E = e V$, where $V$ is the potential energy of the anode, $m$ is the mass of the electron, while $v$ is the speed of the electron after being shot out.

\hfil

\subsection{Exercise 13}

To get the charge-to-mass ratio $\frac{e}{m}$, the initial equation implies the following:
\begin{align}
    e V = \frac{1}{2}mv^2 \implies \frac{e}{m} = \frac{v^2}{2V}
\end{align}

\hfil

\subsection{Exercise 14}

Given the equation derived in \textbf{Exercise 13}, given $\frac{e}{m} = \frac{v^2}{2V}$, we have:
\begin{itemize}
    \item The variable $V$ can be measured directly. Since this represents the anode potential, it is measured using the anode power supply, which is a 2-Channel DC Power Supply in this lab.
    \item The variable $v$ (which is the electron speed) must be deduced using other measurements.
\end{itemize}

\newpage

\section{Magnetic Field Strength}
\subsection{Exercise 15}
Here is the expression of the magnectic field strength:
\begin{equation}
    B = \frac{8}{\sqrt{125}}\frac{\mu_0 N}{a}I
\end{equation}
Which, the given lab apparatus has $\mu_0 = 4\pi \cdot 10^{-7}$ T$\cdot$m/A (permeability of free space),$N = 62\pm 0.1$ (as the number' of turns in the coil), and $a = 33\pm 1$ cm (the radius of the coil). And, $I$ is the curent flowing in the coil.

\hfil

To calculate the propogation of uncertainty in $B$, first it needs the right unit: instead of in cm, take $a = 0.33 \pm 0.01$ m will be the right unit to express $B$ in unit of T.

Then, consider the variables with uncertainty, one has the following partial derivatives:
\begin{align}
    & \frac{\partial B}{\partial N} = \frac{8}{\sqrt{125}}\frac{\mu_0}{a}I,\quad\quad \frac{\partial B}{\partial a} = -\frac{8}{\sqrt{125}}\frac{\mu_0 N}{a^2}I,\quad\quad \frac{\partial B}{\partial I}  = \frac{8}{\sqrt{125}}\frac{\mu_0 N}{a}
\end{align}
Which, let $\delta N,\ \delta a,\ \delta I$ be the uncertainty of number of coil turns, coil radius, and the current, one has the following propogation of uncertainty:
\begin{align}
    \delta B &= \sqrt{\left(\frac{\partial B}{\partial N}\cdot \delta N\right)^2+\left(\frac{\partial B}{\partial a}\cdot \delta a\right)^2+\left(\frac{\partial B}{\partial I}\cdot \delta I\right)^2}\\
    &= \sqrt{\left(\frac{8\mu_0}{\sqrt{125}}\cdot \frac{I}{a}\cdot \delta N\right)^2 + \left(\frac{8\mu_0}{\sqrt{125}}\cdot\frac{NI}{a^2}\cdot\delta a\right)^2 + \left(\frac{8\mu_0}{\sqrt{125}}\cdot\frac{N}{a}\cdot\delta I\right)^2}
\end{align}

\hfil

\subsection{Exercise 16}

Given that $\mu_0 = 4\pi\cdot 10^{-7}$ T$\cdot$m/A, $a = 0.33 \pm 0.01$ m, and $N = 62\pm 0.1$; this shows that $\delta a = 0.01$ m, and $\delta N = 0.1$. On the other hand, the coil has its current provided by the Single Channel DC Power Supply, which has output current $I \in [0,5]$ (unit: A), and uncertainty $\delta I = \in [0.006,0.011]$ A (this is what is measured).

Then, plugin to the three values for $\delta B$, namely $\frac{8\mu_0}{\sqrt{125}}\cdot\frac{I}{a}\cdot \delta N$, $\frac{8\mu_0}{\sqrt{125}}\cdot\frac{N I}{a}\cdot\delta a$, and $\frac{8\mu_0}{\sqrt{125}}\cdot\frac{N}{a}\cdot\delta I$, one has the following:
\begin{align}
    0\leq \frac{8\mu_0}{\sqrt{125}}\cdot \frac{I}{a}\cdot \delta N \leq \frac{8\cdot 4\pi\cdot 10^{-7}}{\sqrt{125}}\cdot\frac{5}{0.33}\cdot 0.1 \approx 1.362\cdot 10^{-6}
\end{align}
\begin{align}
    0\leq \frac{8\mu_0}{\sqrt{125}}\cdot\frac{N I}{a^2}\cdot \delta a \leq \frac{8\cdot 4\pi\cdot 10^{-7}}{\sqrt{125}}\cdot\frac{62\cdot 5}{(0.33)^2}\cdot 0.01 \approx 2.560 \cdot 10^{-5}
\end{align}
\begin{align}
    \frac{8\mu_0}{\sqrt{125}}\cdot\frac{N}{a}\cdot \delta I = \frac{8\cdot 4\pi\cdot 10^{-7}}{\sqrt{125}}\cdot\frac{62}{0.33}\cdot 0.005 \approx 8.447\cdot 10^{-7}
\end{align}
Unfortunately, the first two rely on the exact value of the current $I$, so it can only be provided with the inequality. But, if consider the maximum of all of them, one has:
\begin{itemize}
    \item The largest term being $\frac{8\mu_0}{\sqrt{125}}\cdot\frac{N I}{a^2}\cdot \delta a $.
    \item The smallest term being $\frac{8\mu_0}{\sqrt{125}}\cdot\frac{N}{a}\cdot \delta I$.
    \item If consider the precision up to 3 decimal places, this implies the difference between the scale must be within $100$. Since for the above three terms, the scale difference is only up to $10^2 = 100$, there isn't any term that's really negligible.
    
    (Note: unless the current $I$ is small enough, but within the tested currents, they're all above $1$ A, so this is also not something needed to be considered).
\end{itemize}

\hfil

\section{Circular Motion}
\subsection{Exercise 17}

Given the magnetic force magnitude as follow:
\begin{equation}
    F = evB
\end{equation}
The value $F$ is the magnetic force magnitude, $e$ again represents the electron charge, $v$ also again represents the electron speed after being shot out, while $B$ is the magnetic field strength.

\hfil

\subsection{Exercise 18}

Given the the electrons are going in circular motion (since it can be approximated with uniform speed, and the magnetic field can also be approximated with uniform strength and direction in the setup), then the force on the electron is also the centripetal force. Which, it's given as the follow:
\begin{equation}
    F = \frac{mv^2}{R}
\end{equation}
Here, the value $F$ is the centripetal force magnitude, $m$ is again the electron mass, $v$ also again represents the electron speed after being shot out, while $R$ is the radius of the circular trajectory of the electrons.

\newpage

\section{Model}

\subsection{Exercise 19}

If given that thetwo equations in \textbf{Exercise 17} and \textbf{18} are equal according to the setup, one gets the following:
\begin{align}
    evB = \frac{mv^2}{R}\implies v = RB \frac{e}{m}\implies v^2 = R^2B^2 \left(\frac{e}{m}\right)^2
\end{align}

\hfil

\subsection{Exercise 20}

Recall that initially one had derived the equation $\frac{e}{m} = \frac{v^2}{2V}$ (given $V$ as the electric potential of the anode). Which, given that $v^2 = R^2B^2\left(\frac{e}{m}\right)^2$ derived in \textbf{Exercise 19}, one has the following for $\frac{e}{m}$:
\begin{align}
    \frac{e}{m} = \frac{v^2}{2V} = \frac{R^2B^2}{2V}\left(\frac{e}{m}\right)^2 \implies \frac{e}{m} = \frac{2V}{R^2B^2}
\end{align}

\hfil

\subsection{Exercise 21}
Given the charge to mass ratio $S := \frac{e}{m} = \frac{2V}{R^2B^2}$, since $V,R,B$ are all variables, then they have the corresponding partial derivatives:
\begin{align}
    \frac{\partial S}{\partial V} = \frac{2}{R^2B^2},\quad \quad \frac{\partial S}{\partial R} = -\frac{4V}{R^3B^2},\quad \quad \frac{\partial S}{\partial B} = -\frac{4V}{R^2B^3}
\end{align}
Hence, let $\delta V,\delta R,\delta B$ represent the uncertainty in anode electric potential, electron trajectory radius, and magnetic field strength, the propogation of uncertainty of $S=\frac{e}{m}$ is given as:
\begin{align}
    \delta S &= \sqrt{\left(\frac{\partial S}{\partial V}\cdot \delta V\right)^2+\left(\frac{\partial S}{\partial R}\cdot \delta R\right)^2+\left(\frac{\partial S}{\partial B}\cdot \delta B\right)^2}\\
    &= \sqrt{\left(\frac{2}{R^2B^2}\cdot\delta V\right)^2+\left(\frac{4V}{R^3B^2}\cdot \delta R\right)^2+\left(\frac{4V}{R^2B^3}\cdot\delta B\right)^2}
\end{align}
In particular, one has $\delta V = 0.05$ V (due to limitation of devices' display up to $10$th decimal places), and $\delta R = 0.5$ mm (the measurement of the distance of the cross-bars are in integer multiples of milimeters, which didn't cover the rational part).

\newpage

\section{Analysis}

\subsection{Exercise 22}
Here is the table for experimental data (where $V$ is the accelerating voltage, $D$ is the cross-bar position / distance, $I$ is the average Helmholtz coil current, $R$ is the beam path radius, and $B$ is the magnetic field strength of the coil):
\begin{table}[ht]
\centering
\small
\begin{tabular}{cccccccccc}
\toprule
$V$ (V) & $\delta V$ (V) & $D$ (m) & $\delta D$ (m) & $I$ (A) & $\delta I$ (A) & $R$ (m) & $\delta R$ (m) & $B$ (T) & $\delta B$ (T) \\
\midrule
25.0 & 0.05 & 0.064 & 0.0005 & 2.390 & 0.006 & 0.032 & 0.00025 & $4.038\times10^{-4}$ & $1.229\times10^{-5}$ \\
     &      & 0.078 &        & 2.000 & 0.006 & 0.039 &          & $3.379\times10^{-4}$ & $1.030\times10^{-5}$ \\
     &      & 0.090 &        & 1.710 & 0.010 & 0.045 &          & $2.889\times10^{-4}$ & $8.928\times10^{-6}$ \\
     &      & 0.104 &        & 1.520 & 0.006 & 0.052 &          & $2.568\times10^{-4}$ & $7.858\times10^{-6}$ \\
     &      & 0.114 &        & 1.370 & 0.010 & 0.057 &          & $2.314\times10^{-4}$ & $7.224\times10^{-6}$ \\
\midrule
30.1 & 0.05 & 0.064 & 0.0005 & 2.600 & 0.010 & 0.032 & 0.00025 & $4.392\times10^{-4}$ & $1.344\times10^{-5}$ \\
     &      & 0.078 &        & 2.170 & 0.006 & 0.039 &          & $3.666\times10^{-4}$ & $1.117\times10^{-5}$ \\
     &      & 0.090 &        & 1.870 & 0.011 & 0.045 &          & $3.159\times10^{-4}$ & $9.765\times10^{-6}$ \\
     &      & 0.104 &        & 1.620 & 0.011 & 0.052 &          & $2.737\times10^{-4}$ & $8.510\times10^{-6}$ \\
     &      & 0.114 &        & 1.440 & 0.006 & 0.057 &          & $2.433\times10^{-4}$ & $7.451\times10^{-6}$ \\
\midrule
35.0 & 0.05 & 0.064 & 0.0005 & 2.830 & 0.006 & 0.032 & 0.00025 & $4.781\times10^{-4}$ & $1.454\times10^{-5}$ \\
     &      & 0.078 &        & 2.360 & 0.006 & 0.039 &          & $3.987\times10^{-4}$ & $1.214\times10^{-5}$ \\
     &      & 0.090 &        & 2.020 & 0.010 & 0.045 &          & $3.413\times10^{-4}$ & $1.049\times10^{-5}$ \\
     &      & 0.104 &        & 1.770 & 0.006 & 0.052 &          & $2.990\times10^{-4}$ & $9.130\times10^{-6}$ \\
     &      & 0.114 &        & 1.580 & 0.011 & 0.057 &          & $2.669\times10^{-4}$ & $8.310\times10^{-6}$ \\
\midrule
40.0 & 0.05 & 0.064 & 0.0005 & 3.020 & 0.011 & 0.032 & 0.00025 & $5.102\times10^{-4}$ & $1.559\times10^{-5}$ \\
     &      & 0.078 &        & 2.540 & 0.006 & 0.039 &          & $4.291\times10^{-4}$ & $1.306\times10^{-5}$ \\
     &      & 0.090 &        & 2.170 & 0.006 & 0.045 &          & $3.666\times10^{-4}$ & $1.117\times10^{-5}$ \\
     &      & 0.104 &        & 1.900 & 0.010 & 0.052 &          & $3.210\times10^{-4}$ & $9.886\times10^{-6}$ \\
     &      & 0.114 &        & 1.700 & 0.006 & 0.057 &          & $2.872\times10^{-4}$ & $8.774\times10^{-6}$ \\
\midrule
45.0 & 0.05 & 0.064 & 0.0005 & 3.230 & 0.010 & 0.032 & 0.00025 & $5.457\times10^{-4}$ & $1.664\times10^{-5}$ \\
     &      & 0.078 &        & 2.690 & 0.011 & 0.039 &          & $4.544\times10^{-4}$ & $1.391\times10^{-5}$ \\
     &      & 0.090 &        & 2.330 & 0.011 & 0.045 &          & $3.936\times10^{-4}$ & $1.209\times10^{-5}$ \\
     &      & 0.104 &        & 2.030 & 0.006 & 0.052 &          & $3.429\times10^{-4}$ & $1.046\times10^{-5}$ \\
     &      & 0.114 &        & 1.800 & 0.010 & 0.057 &          & $3.041\times10^{-4}$ & $9.381\times10^{-6}$ \\
\bottomrule
\end{tabular}
\caption{The Rate Data of the Lab}
\end{table}

Here, since the radius $R = \frac{D}{2}$ (since the cross-bar position is the distance between the anode - where the electrons started at, and the cross-bar - where the circular path touches, which represents the diameter of the given circular path). Hence, this implies that $\delta R = \sqrt{\left(\frac{\partial R}{\partial D}\cdot \delta D\right)^2} = \frac{\delta D}{2}$.

\hfil

\subsection{Exercise 23}

Because of the size of the document, the columns for $D$ and $\delta D$ will be removed here. If including $S = \frac{e}{m}$ and $\delta S$ calculated in each case, here is the table:
\begin{table}[ht]
\centering
\small
\begin{tabular}{c|c||c|c||c|c||c|c||c|c}
\toprule
$V$ (V) & $\delta V$ (V)  & $I$ (A) & $\delta I$ (A) &
$R$ (m) & $\delta R$ (m) & $B$ (T) & $\delta B$ (T) & $S$ (C/kg) & $\delta S$ (C/kg) \\
\midrule
\midrule
25.0 & 0.05 & 2.390 & 0.006 & 0.032 & 0.00025 & $4.038\times10^{-4}$ & $1.229\times10^{-5}$ & $2.995\times10^{11}$ & $1.825\times10^{10}$ \\
     &      & 2.000 & 0.006 & 0.039 &          & $3.379\times10^{-4}$ & $1.030\times10^{-5}$ & $2.880\times10^{11}$ & $1.757\times10^{10}$ \\
     &      & 1.710 & 0.010 & 0.045 &          & $2.889\times10^{-4}$ & $8.928\times10^{-6}$ & $2.959\times10^{11}$ & $1.830\times10^{10}$ \\
     &      & 1.520 & 0.006 & 0.052 &          & $2.568\times10^{-4}$ & $7.858\times10^{-6}$ & $2.804\times10^{11}$ & $1.717\times10^{10}$ \\
     &      & 1.370 & 0.010 & 0.057 &          & $2.314\times10^{-4}$ & $7.224\times10^{-6}$ & $2.873\times10^{11}$ & $1.794\times10^{10}$ \\
\midrule
30.1 & 0.05 & 2.600 & 0.010 & 0.032 & 0.00025 & $4.392\times10^{-4}$ & $1.344\times10^{-5}$ & $3.047\times10^{11}$ & $1.865\times10^{10}$ \\
     &      & 2.170 & 0.006 & 0.039 &          & $3.666\times10^{-4}$ & $1.117\times10^{-5}$ & $2.945\times10^{11}$ & $1.796\times10^{10}$ \\
     &      & 1.870 & 0.011 & 0.045 &          & $3.159\times10^{-4}$ & $9.765\times10^{-6}$ & $2.979\times10^{11}$ & $1.842\times10^{10}$ \\
     &      & 1.620 & 0.011 & 0.052 &          & $2.737\times10^{-4}$ & $8.510\times10^{-6}$ & $2.972\times10^{11}$ & $1.849\times10^{10}$ \\
     &      & 1.440 & 0.006 & 0.057 &          & $2.433\times10^{-4}$ & $7.451\times10^{-6}$ & $3.131\times10^{11}$ & $1.919\times10^{10}$ \\
\midrule
35.0 & 0.05 & 2.830 & 0.006 & 0.032 & 0.00025 & $4.781\times10^{-4}$ & $1.454\times10^{-5}$ & $2.991\times10^{11}$ & $1.820\times10^{10}$ \\
     &      & 2.360 & 0.006 & 0.039 &          & $3.987\times10^{-4}$ & $1.214\times10^{-5}$ & $2.895\times10^{11}$ & $1.764\times10^{10}$ \\
     &      & 2.020 & 0.010 & 0.045 &          & $3.413\times10^{-4}$ & $1.049\times10^{-5}$ & $2.968\times10^{11}$ & $1.826\times10^{10}$ \\
     &      & 1.770 & 0.006 & 0.052 &          & $2.990\times10^{-4}$ & $9.130\times10^{-6}$ & $2.895\times10^{11}$ & $1.769\times10^{10}$ \\
     &      & 1.580 & 0.011 & 0.057 &          & $2.669\times10^{-4}$ & $8.310\times10^{-6}$ & $3.024\times10^{11}$ & $1.884\times10^{10}$ \\
\midrule
40.0 & 0.05 & 3.020 & 0.011 & 0.032 & 0.00025 & $5.102\times10^{-4}$ & $1.559\times10^{-5}$ & $3.001\times10^{11}$ & $1.835\times10^{10}$ \\
     &      & 2.540 & 0.006 & 0.039 &          & $4.291\times10^{-4}$ & $1.306\times10^{-5}$ & $2.857\times10^{11}$ & $1.739\times10^{10}$ \\
     &      & 2.170 & 0.006 & 0.045 &          & $3.666\times10^{-4}$ & $1.117\times10^{-5}$ & $2.940\times10^{11}$ & $1.792\times10^{10}$ \\
     &      & 1.900 & 0.010 & 0.052 &          & $3.210\times10^{-4}$ & $9.886\times10^{-6}$ & $2.872\times10^{11}$ & $1.769\times10^{10}$ \\
     &      & 1.700 & 0.006 & 0.057 &          & $2.872\times10^{-4}$ & $8.774\times10^{-6}$ & $2.985\times10^{11}$ & $1.824\times10^{10}$ \\
\midrule
45.0 & 0.05 & 3.230 & 0.010 & 0.032 & 0.00025 & $5.457\times10^{-4}$ & $1.664\times10^{-5}$ & $2.952\times10^{11}$ & $1.801\times10^{10}$ \\
     &      & 2.690 & 0.011 & 0.039 &          & $4.544\times10^{-4}$ & $1.391\times10^{-5}$ & $2.865\times10^{11}$ & $1.755\times10^{10}$ \\
     &      & 2.330 & 0.011 & 0.045 &          & $3.936\times10^{-4}$ & $1.209\times10^{-5}$ & $2.869\times10^{11}$ & $1.762\times10^{10}$ \\
     &      & 2.030 & 0.006 & 0.052 &          & $3.429\times10^{-4}$ & $1.046\times10^{-5}$ & $2.830\times10^{11}$ & $1.726\times10^{10}$ \\
     &      & 1.800 & 0.010 & 0.057 &          & $3.041\times10^{-4}$ & $9.381\times10^{-6}$ & $2.996\times10^{11}$ & $1.849\times10^{10}$ \\
\bottomrule
\end{tabular}
\caption{Tables including Charge to Mass Ratio of the Elecron}
\end{table}

\hfil

\subsection{Exercise 24}

The weighted average based on the above data, is given as follow:
\begin{align}
    \left<\frac{e}{m}\right> = \frac{\sum_{i}\frac{(e/m)_i}{\sigma^2_i}}{\sum_{i}\frac{1}{\sigma^2_i}}= 2.937\cdot 10^{11} \text{ C/kg}
\end{align}

\hfil

\subsection{Exercise 25}

Given the equation $S=\frac{e}{m}=\frac{2V}{R^2B^2}$, if consider the chart with horizontal axis $R^2B^2$, while vertical axis $2V$, then the linear fit of the data points would produce a slop;e of $\frac{e}{m}$. Which, here is the formed charts:

\begin{figure}[h!]
    \centering
    \includegraphics[width=100mm]{scatter_plot.png}
\end{figure}

\newpage

\hfil

\subsection{Exercise 26}

First, if consider the scales of $\frac{\delta S}{S}$ (given that $S:= \frac{e}{m}$), one has the maximum being $0.06245$, while the minimum being $0.06086$. Hence, the uncertainty is within $6.086\%$ to $6.245\%$ in comparison to the derived constant.

This implies the data is neither poor nor excellent (since for poor data, the uncertainty is $>20\%$, while excellent data has $<1\%$ uncertainty).

\hfil

Now, if consider the weighted charge-to-mass ratio, one has $\left<\frac{e}{m}\right> \approx 2.937 \cdot 10^{11}$ C/kg. Going on to the National Institute of Science and Technology's website about Fundamental Physical Constants, the generally accepted value is $\approx 1.759 \cdot 10^{11}$ C/kg.

Which, the discrepency (or percent error) is about $66.970\%$, which is a lot greater than the maximum uncertainty within the data.

(Source citation: 

``The NIST Reference on Constants, Units, and Uncertainty.'' 2022, 

https://physics.nist.gov/cuu/Constants/bibliography.html ).

\hfil

To talk about some potential sources of error, here are some possibilities:
\begin{itemize}
    \item During the start of the derivation, there is an equation $eV = \frac{1}{2}mv^2$, and such equation is based on the assumption that initially, the electrons before bing shot out (and before the anode was provided any potential) have no energy at all. However, there's no measurement one can do in this experiment to verify that, and there's a possibility that the initial electrons in fact held some nonzero energy, and causing the experimental value to differentiate from the expected values.
\end{itemize}
Statistical wise, the largest uncertainty should be the fact that there's fewer measurements regarding all quantities. In this experiment (as a single person group), I spent the majority of time on constructing the equipments and wiring them, which left with only less than an hour to collect all the required data. This unfortunately provides little to no time to repeat the experiments for multiple trials regarding every variables, so the only estimated uncertainty is done on the coil current, and everything else (for instance, the voltage) has uncertainty being the limitation the devices had, instead of actually being estimated. 

This would provide extra uncertainties within the data (since there are fewer data set to do comparison with), and is a likely source of error. Hence, in the future to decrease such uncertainty, being more efficient during the initial preparation and leave more time for data collection is necessary.

\hfil

\section{Conclusion (Exercise 27)}

Through this week's experiment, with the collected data on the anode voltage (which corresponds to the electrons' kinetic energy), the Helmholtz coil's current (which corresponds to the magnetic field strength), and the cross-bar distance (which corresponds to the radius of the electrons' trajectory), the data were used to calculate the electron Charge-to-Mass ratio. 

\hfil

Here, the collected data derived a weighted average of the ratio to be $\left<\frac{e}{m}\right> = 2.937\cdot 10^{11}$ C/kg, together with all the uncertainties to be in the range of $6.086\%$ to $6.245\%$, showing the data is neither excellent nor poor.

However, in comparison to the generaly accepted value of $1.759\cdot 10^{11}$ C/kg, the weighted average provides a percent error of $66.970\%$, which is a lot larger in comparison to the uncertainty derived in this experiment.

\hfil

Talking about the sources of error:

First, the initial assumption that the free electrons carry no energy leads to the equation $e V=\frac{1}{2}mv^2$ used in all the derivations, which may be a false assumption. Yet, it's not possible to verify such claim, at least not within the experimental setup.

Second, because of the time constraint there is less time for the data collection, which the experiment can't be performed with multiple trials regarding each independent variables. This could bring extra statistical uncertainties to the data, and causing the huge difference between the experimental result with the generally accepted values. To fix this problem, being more efficient in the initial setup and preserve more time for data collection is necessary.

\end{document}